\section*{Chapitre \ref{ch:b2:integration} --- Intégration}

\begin{solution}{exo:b2:integration:1}
$\int_0^1 \frac{\dd t}{\sqrt{t(1-t)}}$~: au voisinage de $0$, $\sim t^{-1/2}$
($\alpha = \frac12 < 1$~: \omterm{def:b2:integration:improper}{converge})~; au voisinage de $1$, $\sim
(1-t)^{-1/2}$~: \omterm{def:b2:integration:improper}{converge}. Convergente (sa valeur est $\pi$, par le
changement de variable $t = \sin^2\theta$).

$\int_1^\infty \frac{\ln t}{t^2}$~: $\frac{\ln t}{t^2} = o(t^{-3/2})$~:
convergente (valeur $1$ par parties).

$\int_0^\infty \frac{\dd t}{1 + t^2\sin^2 t}$~: divergente. Au voisinage de $t =
n\pi$, écrivons $t = n\pi + u$~: $\sin^2 t = \sin^2 u \leq u^2$, donc sur
$\abs u \leq \frac{1}{n}$, $1 + t^2\sin^2 t \leq 1 + (n\pi +
1)^2u^2 \leq C n^2 u^2 + 1$~; d'où
\[
\int_{n\pi - 1/n}^{n\pi + 1/n} \frac{\dd t}{1 + t^2\sin^2 t}
\geq \int_{-1/n}^{1/n} \frac{\dd u}{1 + Cn^2u^2}
= \frac{2\arctan\sqrt C}{\sqrt C}\cdot\frac{1}{n} ,
\]
un terme d'une série divergente de type harmonique~: en sommant sur $n$, la
primitive est non bornée.
\end{solution}

\begin{solution}{exo:b2:integration:2}
Substituons $u = \lambda t$~:
\[
\int_0^\infty t^n \eu^{-\lambda t}\dd t
= \frac{1}{\lambda^{n+1}}\int_0^\infty u^n\eu^{-u}\dd u
= \frac{\Gamma(n+1)}{\lambda^{n+1}} = \frac{n!}{\lambda^{n+1}} .
\]
Avec $t = \eu^{-u}$ ($\dd t = -\eu^{-u}\dd u$)~:
\[
\int_0^1 (\ln t)^n \dd t = \int_0^{\infty} (-u)^n \eu^{-u}\,\dd u
= (-1)^n\, n! .
\]
\end{solution}

\begin{solution}{exo:b2:integration:3}
Convergence~: l'intégrande est $\leq \frac{1}{1+t^2}$ au voisinage de $\infty$
et borné au voisinage de $0$ (les deux facteurs minorés loin de $0$)~:
\omterm{def:b2:series:def}{absolument} convergente, pour tout $x$. En substituant $t = \frac1u$
($\dd t = -\frac{\dd u}{u^2}$)~:
\[
I(x) = \int_0^\infty \frac{1}{\bigl(1 +
\frac1{u^2}\bigr)\bigl(1 + u^{-x}\bigr)}\cdot\frac{\dd u}{u^2}
= \int_0^\infty \frac{u^x}{(1 + u^2)(1 + u^x)}\,\dd u .
\]
En additionnant les deux expressions de $I(x)$~:
\[
2I(x) = \int_0^\infty \frac{1 + t^x}{(1+t^2)(1+t^x)}\dd t
= \int_0^\infty \frac{\dd t}{1 + t^2} = \frac{\pi}{2} :
\]
$I(x) = \frac\pi4$, indépendant de $x$.
\end{solution}

\begin{solution}{exo:b2:integration:4}
Au voisinage de $0^+$, avec $u = \abs{\ln t} \to \infty$. Si $\alpha < 1$~:
convergence quel que soit $\beta$ (comparer avec $t^{-\alpha'}$ pour
$\alpha < \alpha' < 1$~: le facteur logarithmique est battu). Si $\alpha > 1$~:
divergence quel que soit $\beta$ (comparer avec $t^{-\alpha''}$,
$1 < \alpha'' < \alpha$). Si $\alpha = 1$~: substituer $t =
\eu^{-u}$~:
\[
\int_0^{1/2} \frac{\dd t}{t\,\abs{\ln t}^\beta}
= \int_{\ln 2}^{\infty} \frac{\dd u}{u^\beta},
\]
convergente si et seulement si $\beta > 1$. Récapitulatif~: convergence si et seulement si $\alpha < 1$,
ou ($\alpha = 1$ et $\beta > 1$) --- le miroir de la série de Bertrand.
\end{solution}

\begin{solution}{exo:b2:integration:5}
Continuité sur $\intco{0}{\infty}$~: domination
$\bigl|\frac{\eu^{-xt}}{1+t^2}\bigr| \leq \frac{1}{1+t^2}$,
intégrable, uniforme en $x \geq 0$~:
le~\cref{thm:b2:integration:continuity}.

$C^2$ sur $\intoo{0}{\infty}$~: sur $x \geq a > 0$, les deux premières
dérivées en $x$, $\frac{-t\,\eu^{-xt}}{1+t^2}$ et
$\frac{t^2\eu^{-xt}}{1+t^2}$, sont dominées par $t\,\eu^{-at}$ et
$\eu^{-at}$~: deux applications du~\cref{thm:b2:integration:leibnizrule}.
Puis
\[
F''(x) + F(x) = \int_0^\infty \frac{t^2 + 1}{1 + t^2}\,\eu^{-xt}\dd
t = \int_0^\infty \eu^{-xt}\dd t = \frac1x .
\]
Limite~: $0 \leq F(x) \leq \int_0^\infty \eu^{-xt}\dd t = \frac1x \to
0$.
\end{solution}

\begin{solution}{exo:b2:integration:6}
Sur $\intcc{\varepsilon}{M}$, substituons $u = at$ et $u = bt$ dans
les deux moitiés~:
\[
\int_\varepsilon^M \frac{f(at) - f(bt)}{t}\dd t
= \int_{a\varepsilon}^{aM}\frac{f(u)}{u}\dd u
- \int_{b\varepsilon}^{bM}\frac{f(u)}{u}\dd u
= \int_{a\varepsilon}^{b\varepsilon} \frac{f(u)}{u}\dd u
- \int_{aM}^{bM} \frac{f(u)}{u}\dd u .
\]
Premier morceau~: $f(u) = f(0) + o(1)$ au voisinage de $0$, et $\int_{a\varepsilon}
^{b\varepsilon} \frac{\dd u}{u} = \ln\frac ba$~: le morceau tend vers
$f(0)\ln\frac ba$. Second morceau~: $f(u) \to f(\infty)$, même
calcul~: tend vers $f(\infty)\ln\frac ba$. D'où l'\omterm{def:b2:integration:improper}{intégrale impropre}
\omterm{def:b2:integration:improper}{converge} vers $\bigl(f(0) - f(\infty)\bigr)\ln\frac ba$.

Avec $f(t) = \eu^{-t}$ ($f(0) = 1$, $f(\infty) = 0$), $a = 1$, $b
= 2$~:
\[
\int_0^\infty \frac{\eu^{-t} - \eu^{-2t}}{t}\dd t = \ln 2 .
\]
\end{solution}

\begin{solution}{exo:b2:integration:7}
Prolongeons l'intégrande par $0$ au-delà de $t = n$~: $g_n(t) = (1 -
\frac tn)^n t^{x-1}\mathbf{1}_{t \leq n}$. Simplement, $g_n(t) \to
\eu^{-t}t^{x-1}$ (la limite des intérêts composés, volume de première année).
Domination~: $\ln(1 - u) \leq -u$ donne $(1 - \frac tn)^n \leq
\eu^{-t}$ sur $\intcc{0}{n}$, donc $\abs{g_n(t)} \leq
\eu^{-t}t^{x-1} = \varphi(t)$, intégrable. Convergence dominée~:
\[
\int_0^n \Bigl(1 - \frac tn\Bigr)^n t^{x-1}\dd t
\xrightarrow[n\to\infty]{} \int_0^\infty \eu^{-t}t^{x-1}\dd t
= \Gamma(x) .
\]
(Le calcul du membre de gauche par parties répétées donne la forme
produit d'Euler $\Gamma(x) = \lim \frac{n!\,n^x}{x(x+1)\cdots(x+n)}$.)
\end{solution}

\begin{solution}{exo:b2:integration:8}
$H$ est dérivable en $x$ (intégrande $C^1$ en $x$, dérivée
$-2x(1+t^2)\cdot\frac{\eu^{-x^2(1+t^2)}}{1+t^2} =
-2x\,\eu^{-x^2}\eu^{-x^2t^2}$, \omterm{def:b2:metric:continuity}{continue} et bornée sur les \omterm{def:b2:metric:compact}{compacts}
de $x$, domination sur $t \in \intcc{0}{1}$ triviale)~:
\[
H'(x) = -2x\,\eu^{-x^2}\int_0^1 \eu^{-x^2t^2}\,\dd t
\overset{u = xt}{=} -2\,\eu^{-x^2}\int_0^x \eu^{-u^2}\,\dd u
= -G'(x),
\]
puisque $G'(x) = 2\eu^{-x^2}\int_0^x \eu^{-t^2}\dd t$ (règle de la chaîne sur
le carré, théorème fondamental de l'analyse). Donc $G + H$ est
constante, égale à $G(0) + H(0) = 0 + \int_0^1 \frac{\dd t}{1+t^2}
= \frac\pi4$.

Quand $x \to \infty$~: $0 \leq H(x) \leq \eu^{-x^2}\int_0^1 \dd t \to
0$, donc $G(x) \to \frac\pi4$~:
\[
\int_0^\infty \eu^{-t^2}\dd t = \sqrt{\frac\pi4} =
\frac{\sqrt\pi}{2} .
\]
(Par conséquent $\Gamma\bigl(\frac12\bigr) = 2\int_0^\infty
\eu^{-t^2}\dd t = \sqrt\pi$, par le changement de variable $t = \sqrt u$.)
\end{solution}

\begin{solution}{exo:b2:integration:9}
Cauchy--Schwarz (volume de première année, valable sur $\intcc{\varepsilon}{M}$
et passée à la limite) appliquée à la factorisation
$t^{\frac{x+y}{2}-1}\eu^{-t} = \bigl(t^{x-1}\eu^{-t}\bigr)^{1/2}
\bigl(t^{y-1}\eu^{-t}\bigr)^{1/2}$~:
\[
\Gamma\Bigl(\frac{x+y}{2}\Bigr) \leq
\Gamma(x)^{1/2}\,\Gamma(y)^{1/2}
\quad\Longrightarrow\quad
\ln\Gamma\Bigl(\frac{x+y}{2}\Bigr) \leq \frac{\ln\Gamma(x) +
\ln\Gamma(y)}{2} :
\]
$\ln\Gamma$ est convexe au milieu~; étant \omterm{def:b2:metric:continuity}{continue}
(\cref{thm:b2:integration:gammaprops}), elle est convexe
(\cref{exo:b2:realfun:8}). (La log-convexité caractérise $\Gamma$ de
manière unique parmi les interpolations de la factorielle --- le
théorème de Bohr--Mollerup, une perle de troisième année.)
\end{solution}

\begin{solution}{exo:b2:integration:10}
\begin{enumerate}
  \item Sur $x \geq a > 0$~: la dérivée en $x$ de l'intégrande est
        $-\sin t\,\eu^{-xt}$, dominée par $\eu^{-at}$~:
        le~\cref{thm:b2:integration:leibnizrule} donne $F'(x) =
        -\int_0^\infty \eu^{-xt}\sin t\,\dd t$. Deux intégrations par
        parties (ou l'exponentielle complexe)~:
        \[
        \int_0^\infty \eu^{-xt}\sin t\,\dd t
        = \Im \int_0^\infty \eu^{(-x+\iu)t}\dd t
        = \Im\frac{1}{x - \iu} = \frac{1}{1 + x^2} .
        \]
  \item $\abs{F(x)} \leq \int_0^\infty \eu^{-xt}\dd t = \frac1x \to
        0$. En intégrant $F' = -\frac{1}{1+x^2}$ de $x$ à
        $\infty$~: $0 - F(x) = -\bigl(\frac\pi2 - \arctan x\bigr)$,
        donc $F(x) = \frac\pi2 - \arctan x$.
  \item En faisant $x \to 0^+$ avec la continuité admise~: $F(0^+) =
        \frac\pi2$, et $F(0) = \int_0^\infty \frac{\sin t}{t}\dd
        t$ (l'intégrale de Dirichlet semi-convergente,
        \cref{ex:b2:integration:sint})~: sa valeur est
        $\frac\pi2$.
\end{enumerate}
\end{solution}

\begin{solution}{exo:b2:integration:11}
Convergence~: au voisinage de $0$ l'intégrande se prolonge par continuité par la
valeur $1$ ($\sin t \sim t$)~; à l'infini il est $\leq t^{-2}$~:
convergence absolue. Sur $\intcc{\varepsilon}{M}$, intégrons par
parties avec $u = \sin^2 t$, $v' = t^{-2}$~:
\[
\int_\varepsilon^M \frac{\sin^2 t}{t^2}\dd t
= \Bigl[-\frac{\sin^2 t}{t}\Bigr]_\varepsilon^M
+ \int_\varepsilon^M \frac{2\sin t\cos t}{t}\dd t
= \Bigl[-\frac{\sin^2 t}{t}\Bigr]_\varepsilon^M
+ \int_{2\varepsilon}^{2M} \frac{\sin u}{u}\dd u
\]
($u = 2t$ dans la dernière intégrale). Le crochet tend vers $0$ aux deux
bornes ($\sin^2\varepsilon/\varepsilon \leq \varepsilon$~;
$\sin^2 M/M \leq 1/M$), et la dernière intégrale tend vers
$\int_0^\infty \frac{\sin u}{u}\dd u = \frac\pi2$
(\cref{exo:b2:integration:10}). D'où
\[
\int_0^\infty \Bigl(\frac{\sin t}{t}\Bigr)^{\!2}\dd t
= \frac\pi2 .
\]
\end{solution}

\begin{solution}{exo:b2:integration:12}
\begin{enumerate}
  \item Parties avec $u = \frac{-1}{2t}$, $v' = -2t\,\eu^{-t^2}$
        (donc $v = \eu^{-t^2}$)~:
        \[
        T(x) = \Bigl[\frac{-\eu^{-t^2}}{2t}\Bigr]_x^\infty
        - \int_x^\infty \frac{\eu^{-t^2}}{2t^2}\dd t
        = \frac{\eu^{-x^2}}{2x}
        - \int_x^\infty \frac{\eu^{-t^2}}{2t^2}\dd t .
        \]
        Même procédé sur la nouvelle intégrale ($u = \frac{-1}{4t^3}$,
        $v' = -2t\,\eu^{-t^2}$)~:
        \[
        \int_x^\infty \frac{\eu^{-t^2}}{2t^2}\dd t
        = \frac{\eu^{-x^2}}{4x^3}
        - \frac34\int_x^\infty \frac{\eu^{-t^2}}{t^4}\dd t ,
        \]
        d'où l'identité annoncée.
  \item Une intégration par parties de plus majore le reste~:
        \[
        \int_x^\infty \frac{\eu^{-t^2}}{t^4}\dd t
        = \frac{\eu^{-x^2}}{2x^5}
        - \frac52\int_x^\infty\frac{\eu^{-t^2}}{t^6}\dd t
        \leq \frac{\eu^{-x^2}}{2x^5},
        \]
        donc $0 \leq \frac34\int_x^\infty t^{-4}\eu^{-t^2}\dd t
        \leq \frac{3}{8x^5}\eu^{-x^2}$. En laissant tomber le reste (positif)
        dans l'identité de la question 1 on obtient la borne
        inférieure~; en laissant tomber le deuxième terme (négatif) des premières
        parties on obtient $T(x) \leq \frac{\eu^{-x^2}}{2x}$. En divisant
        l'encadrement par $\frac{\eu^{-x^2}}{2x}$~: le rapport est
        coincé entre $1 - \frac{1}{2x^2}$ et $1$, donc $T(x)
        \sim \frac{\eu^{-x^2}}{2x}$.
  \item En itérant les parties on produit la série formelle
        \[
        T(x) \approx \frac{\eu^{-x^2}}{2x}\Bigl(1 -
        \frac{1}{2x^2} + \frac{1\cdot3}{(2x^2)^2} -
        \frac{1\cdot3\cdot5}{(2x^2)^3} + \cdots\Bigr),
        \]
        dont le coefficient d'\omterm{def:b2:structures:generated}{ordre} $k$, $1\cdot3\cdots(2k-1) =
        \frac{(2k)!}{2^k k!}$, croît plus vite que toute suite
        géométrique~: pour $x$ fixé les termes
        $\frac{1\cdot3\cdots(2k-1)}{(2x^2)^k}$ tendent vers l'infini
        (leur rapport est $\frac{2k+1}{2x^2} \to \infty$), donc la
        série diverge pour tout $x$. C'est un
        développement \emph{asymptotique}~: tronqué à un \omterm{def:b2:structures:generated}{ordre}
        fixé quelconque, l'erreur est de l'\omterm{def:b2:structures:generated}{ordre} du premier terme
        omis quand $x \to \infty$ --- mais jamais une série
        convergente. (Cette estimation de queue est la borne de queue
        gaussienne standard des chapitres de probabilités.)
\end{enumerate}
\end{solution}

\begin{solution}{pb:b2:integration:1}
\textbf{1.} En $0^+$ l'intégrande est $\sim t^{x-1}$~: l'échelle
en une borne finie \omterm{def:b2:integration:improper}{converge} si et seulement si $1 - x < 1$, c.-à-d.\ $x > 0$
(et pour $x \leq 0$, $t^{x-1} \geq t^{-1}$ diverge)~; en
$+\infty$, $t^{x-1}\eu^{-t} = o(t^{-2})$ \omterm{def:b2:integration:improper}{converge} pour tout
$x$. Puis $\Gamma(x) = \frac{\Gamma(x+1)}{x}$ et $\Gamma(x+1)
\to \Gamma(1) = 1$ quand $x \to 0^+$ (continuité,
\cref{thm:b2:integration:gammaprops})~: $\Gamma(x) \sim
\frac1x$.

\textbf{2.} Avec $t = u^2$, $\dd t = 2u\,\dd u$~:
\[
\Gamma\Bigl(\frac12\Bigr) = \int_0^\infty
t^{-1/2}\eu^{-t}\dd t
= \int_0^\infty \frac{\eu^{-u^2}}{u}\,2u\,\dd u
= 2\int_0^\infty \eu^{-u^2}\dd u = \sqrt\pi
\]
d'après l'\cref{exo:b2:integration:8}. Avec $u = v/\sqrt2$~:
\[
\int_\R \eu^{-v^2/2}\dd v
= 2\sqrt2\int_0^\infty \eu^{-u^2}\dd u
= \sqrt2\,\sqrt\pi = \sqrt{2\pi} .
\]

\textbf{3.} Vrai pour $n = 0$ (les deux membres valent $\sqrt\pi$). Si
$\Gamma(n + \frac12) = \frac{(2n)!}{4^n n!}\sqrt\pi$, l'équation
fonctionnelle donne
\[
\Gamma\Bigl(n + 1 + \frac12\Bigr)
= \Bigl(n + \frac12\Bigr)\Gamma\Bigl(n + \frac12\Bigr)
= \frac{2n+1}{2}\cdot\frac{(2n)!}{4^n n!}\sqrt\pi
= \frac{(2n+2)!}{4^{n+1}(n+1)!}\sqrt\pi ,
\]
la dernière étape car $\frac{(2n+2)!}{(2n)!} = (2n+2)(2n+1)$
et $\frac{2n+1}{2} = \frac{(2n+2)(2n+1)}{4(n+1)}$.

\textbf{4.} Le~\cref{thm:b2:integration:gammaprops} donne
$\Gamma''(x) = \int_0^\infty t^{x-1}\eu^{-t}(\ln t)^2\dd t$
(deux applications de la règle de Leibniz, dominations comme dans la
preuve du théorème)~; l'intégrande est $\geq 0$ et non identiquement
nul, donc $\Gamma'' > 0$~: $\Gamma$ est strictement convexe et
$\Gamma'$ strictement croissante. Puisque $\Gamma(1) = \Gamma(2)
= 1$, Rolle fournit $x_0 \in \intoo12$ tel que $\Gamma'(x_0) =
0$~; la stricte monotonie de $\Gamma'$ fait de $x_0$ son unique
zéro, avec $\Gamma' < 0$ avant et $\Gamma' > 0$ après~:
$\Gamma$ décroît sur $\intoo0{x_0}$, croît sur
$\intoo{x_0}\infty$, et $x_0$ est l'unique minimum.

\textbf{5.} Soit $k \in \N$ et $x \geq 3$~; choisissons l'entier
$n$ tel que $n + 1 \leq x < n + 2$ (donc $n \geq 1$). Par la
monotonie de la question 4 (valable à partir de $x_0 < 2$)~: $\Gamma(x)
\geq \Gamma(n + 1) = n!$, tandis que $x^k \leq (n+2)^k$. Puisque
$\frac{n!}{(n+2)^k} \to \infty$ (les factorielles l'emportent sur les puissances,
volume de première année), $\frac{\Gamma(x)}{x^k} \geq \frac{n!}{(n+2)^k} \to
\infty$ quand $x \to \infty$~: $x^k = o(\Gamma(x))$.

\textbf{6.} Au voisinage de $0$ l'intégrande est $\sim t^{x-1}$
(convergent si et seulement si $x > 0$), au voisinage de $1$ il est $\sim (1-t)^{y-1}$
(si et seulement si $y > 0$)~; les deux comparaisons sont entre fonctions positives,
donc $B(x,y)$ \omterm{def:b2:integration:improper}{converge} exactement pour $x, y > 0$. Le changement de variable
$t \mapsto 1 - t$ échange les deux facteurs~: $B(x,y) = B(y,x)$.

\textbf{7.} $B(x,1) = \int_0^1 t^{x-1}\dd t = \frac1x$. Parties
sur $\intcc\varepsilon{1-\varepsilon}$ avec $u = (1-t)^y$, $v =
\frac{t^x}{x}$~:
\[
\int t^{x-1}(1-t)^{y}\dd t
= \Bigl[\frac{t^x(1-t)^y}{x}\Bigr]
+ \frac{y}{x}\int t^{x}(1-t)^{y-1}\dd t ;
\]
le crochet s'annule aux deux bornes quand $\varepsilon \to 0$ ($x >
0$ en $0$, $y > 0$ en $1$), laissant $B(x, y+1) = \frac
yx\,B(x+1, y)$.

\textbf{8.} Puisque $t + (1-t) = 1$~:
\[
t^{x-1}(1-t)^{y-1} = t^{x}(1-t)^{y-1} + t^{x-1}(1-t)^{y},
\]
donc $B(x,y) = B(x+1,y) + B(x,y+1)$. La question 7 s'écrit $B(x+1,y) =
\frac xy B(x,y+1)$~; en substituant,
\[
B(x,y) = \Bigl(\frac xy + 1\Bigr)B(x,y+1)
= \frac{x+y}{y}\,B(x,y+1),
\]
c.-à-d.\ $B(x,y+1) = \frac{y}{x+y}B(x,y)$~; la relation jumelle
découle de la symétrie de la question 6.

\textbf{9.} Récurrence sur $n$ à $m$ fixé~: $B(m,1) = \frac1m =
\frac{(m-1)!\,0!}{m!}$, et si la formule vaut au rang $n$,
\[
B(m, n+1) = \frac{n}{m+n}\,B(m,n)
= \frac{n}{m+n}\cdot\frac{(m-1)!(n-1)!}{(m+n-1)!}
= \frac{(m-1)!\,n!}{(m+n)!} .
\]
En réécrivant~: $B(m,n) = \frac{(m-1)!(n-1)!}{(m+n-1)!} =
\bigl[(m+n-1)\binom{m+n-2}{m-1}\bigr]^{-1}$.

\textbf{10.} Les deux membres de $B(x,n) =
\frac{\Gamma(x)\Gamma(n)}{\Gamma(x+n)}$ valent $\frac1x$ en $n =
1$ ($\Gamma(1) = 1$, $\Gamma(x+1) = x\Gamma(x)$). S'ils coïncident
au rang $n$, alors par la relation de descente et l'équation
fonctionnelle~:
\[
B(x, n+1) = \frac{n}{x+n}\,B(x,n), \qquad
\frac{\Gamma(x)\Gamma(n+1)}{\Gamma(x+n+1)}
= \frac{n}{x+n}\cdot
\frac{\Gamma(x)\Gamma(n)}{\Gamma(x+n)} :
\]
les deux suites vérifient la même récurrence à partir de la même
donnée initiale, donc coïncident pour tout $n \in \N^*$ et tout $x > 0$.

\textbf{11.} Avec $t = \sin^2\theta$ ($\theta \in
\intoo0{\pi/2}$, $\dd t = 2\sin\theta\cos\theta\,\dd\theta$),
$t^{x-1} = \sin^{2x-2}\theta$ et $(1-t)^{y-1} =
\cos^{2y-2}\theta$~:
\[
B(x,y) = \int_0^{\pi/2}\sin^{2x-2}\theta\,\cos^{2y-2}\theta
\cdot 2\sin\theta\cos\theta\,\dd\theta
= 2\int_0^{\pi/2}\sin^{2x-1}\theta\,\cos^{2y-1}\theta\,
\dd\theta .
\]

\textbf{12.} Prenons $y = \frac12$ (annulant le facteur cosinus) et
$2x - 1 = n$~: $B\bigl(\frac{n+1}2, \frac12\bigr) = 2W_n$, c.-à-d.\
$W_n = \frac12 B\bigl(\frac{n+1}2,\frac12\bigr)$. La relation de descente
dans la première variable donne
\[
\frac{W_n}{W_{n-2}}
= \frac{B\bigl(\frac{n-1}2 + 1, \frac12\bigr)}
{B\bigl(\frac{n-1}2, \frac12\bigr)}
= \frac{\frac{n-1}2}{\frac{n-1}2 + \frac12}
= \frac{n-1}{n} :
\]
la récurrence de Wallis, cette fois sans aucune intégration par parties
sur les sinus --- la partie II a fait le travail une fois pour toutes.

\textbf{13.} $B\bigl(\frac12,\frac12\bigr) = 2W_0 =
2\cdot\frac\pi2 = \pi$, tandis que
$\Gamma\bigl(\frac12\bigr)^2/\Gamma(1) = (\sqrt\pi)^2 = \pi$~:
la formule d'Euler est vérifiée en $\bigl(\frac12,\frac12\bigr)$.

\textbf{14.} En itérant $W_{2n} = \frac{2n-1}{2n}W_{2n-2}$ à partir de
$W_0 = \frac\pi2$~:
\[
W_{2n} = \frac\pi2\prod_{k=1}^{n}\frac{2k-1}{2k}
= \frac\pi2\cdot\frac{(2n)!}{4^n(n!)^2},
\]
puisque $\prod(2k-1) = \frac{(2n)!}{2^n n!}$ et $\prod 2k = 2^n
n!$. D'où, en utilisant la question 3~:
\[
B\Bigl(n+\frac12, \frac12\Bigr) = 2W_{2n}
= \pi\,\frac{(2n)!}{4^n(n!)^2}
= \frac{(2n)!\sqrt\pi}{4^n n!}\cdot\frac{\sqrt\pi}{n!}
= \frac{\Gamma\bigl(n+\frac12\bigr)\Gamma\bigl(\frac12\bigr)}
{\Gamma(n+1)} .
\]
Fixons maintenant $x \in \frac12\N^*$. La formule d'Euler est vérifiée en $(x,
\frac12)$~: pour $x$ entier c'est la question 10 (avec la symétrie),
pour $x = n + \frac12$ c'est la formule ci-dessus. Les deux membres de
la formule d'Euler vérifient la récurrence de descente $y \mapsto y + 1$
(question 8 à gauche, l'équation fonctionnelle à droite,
comme à la question 10)~: la récurrence propage la formule de $y =
\frac12$ et $y = 1$ à tout $y \in \frac12\N^*$. La formule d'Euler
vaut donc dès que $2x, 2y \in \N^*$.

\textbf{15.} Avec $u = \frac{t}{1-t}$, c.-à-d.\ $t =
\frac{u}{1+u}$, $1 - t = \frac{1}{1+u}$, $\dd t =
\frac{\dd u}{(1+u)^2}$~:
\[
B(x,y) = \int_0^\infty
\Bigl(\frac{u}{1+u}\Bigr)^{x-1}
\Bigl(\frac{1}{1+u}\Bigr)^{y-1}
\frac{\dd u}{(1+u)^2}
= \int_0^\infty \frac{u^{x-1}}{(1+u)^{x+y}}\,\dd u .
\]
En $x = y = \frac12$, avec $u = v^2$~:
\[
\int_0^\infty \frac{u^{-1/2}}{1+u}\dd u
= \int_0^\infty \frac{2\,\dd v}{1+v^2} = \pi
= B\Bigl(\frac12,\frac12\Bigr) . \checkmark
\]

\textbf{16.} Une intégration par parties, pour $1 \leq k \leq n$
et $s > 0$ ($u = (1 - t/n)^k$, $v = t^s/s$~; les termes de bord
s'annulent)~:
\[
\int_0^n \Bigl(1-\frac tn\Bigr)^{\!k} t^{s-1}\dd t
= \frac{k}{ns}\int_0^n \Bigl(1-\frac tn\Bigr)^{\!k-1}
t^{s}\dd t .
\]
À partir de $k = n$, $s = x$ et en itérant $n$ fois~:
\[
\int_0^n \Bigl(1-\frac tn\Bigr)^{\!n} t^{x-1}\dd t
= \frac{n(n-1)\cdots1}{n^n\,x(x+1)\cdots(x+n-1)}
\int_0^n t^{x+n-1}\dd t
= \frac{n!}{n^n}\cdot
\frac{n^{x+n}}{x(x+1)\cdots(x+n)} ,
\]
soit $\dfrac{n!\,n^x}{x(x+1)\cdots(x+n)}$.

\textbf{17.} D'après l'\cref{exo:b2:integration:7} le membre de gauche tend
vers $\Gamma(x)$ (convergence dominée avec dominant
$t^{x-1}\eu^{-t}$)~; le membre de droite est le quotient de Gauss~:
\[
\Gamma(x) = \lim_{n\to\infty}
\frac{n!\,n^x}{x(x+1)\cdots(x+n)} .
\]

\textbf{18.} En prenant les logarithmes dans le quotient $G_n(x)$
de la question 16 et en découpant $\ln(x+k) = \ln k + \ln(1 + x/k)$ pour $k
\geq 1$~:
\[
\ln G_n(x) = x\ln n - \ln x - \sum_{k=1}^n
\ln\Bigl(1+\frac xk\Bigr)
= -\ln x + x(\ln n - H_n)
+ \sum_{k=1}^n\Bigl(\frac xk -
\ln\Bigl(1+\frac xk\Bigr)\Bigr).
\]
Pour $u \geq 0$, $u - \frac{u^2}2 \leq \ln(1+u) \leq u$, donc le
terme général appartient à $\intcc{0}{x^2/(2k^2)}$~: la série
\omterm{def:b2:integration:improper}{converge} (comparaison avec $\sum k^{-2}$). Puisque $\ln n - H_n
\to -\gamma$ (\cref{ex:b2:comparison:harmonic}) et $\ln G_n(x)
\to \ln\Gamma(x)$ (question 17 et continuité de $\ln$)~:
\[
\ln\Gamma(x) = -\ln x - \gamma x
+ \sum_{k=1}^\infty\Bigl(\frac xk -
\ln\Bigl(1+\frac xk\Bigr)\Bigr) .
\]

\textbf{19.} En $x = \frac12$, le dénominateur est
$\prod_{k=0}^n\bigl(k+\frac12\bigr) =
\frac{(2n+1)!}{2^{2n+1}n!}$ (en développant les moitiés), donc
\[
G_n\Bigl(\frac12\Bigr)
= \frac{n!\,\sqrt n\;2^{2n+1}n!}{(2n+1)!}
= \frac{2\sqrt n\;4^n}{(2n+1)\binom{2n}{n}} .
\]
Avec $\binom{2n}{n} \sim \frac{4^n}{\sqrt{\pi n}}$
(\cref{ex:b2:comparison:centralbinomial})~:
\[
G_n\Bigl(\frac12\Bigr)
\sim \frac{2\sqrt n\,\sqrt{\pi n}}{2n+1}
\longrightarrow \sqrt\pi = \Gamma\Bigl(\frac12\Bigr) .
\]
Le $\sqrt\pi$ du coefficient binomial central (qui vient
de Wallis, donc de la constante de Stirling) et l'intégrale
de Gauss avec son $\sqrt\pi$ sont le même nombre.

\textbf{20.} L'identité relève d'un calcul algébrique direct~: multiplier
$\frac{n!\,n^x}{x(x+1)\cdots(x+n)}$ par $\frac{nx}{x+n+1}$ et
absorber $x$ dans le produit, $n$ dans $n^x$. En faisant $n \to
\infty$~: le membre de gauche tend vers $\Gamma(x+1)$ (Gauss en $x+1$),
le membre de droite vers $\Gamma(x)\cdot x\cdot 1$ puisque
$\frac{n}{x+n+1} \to 1$~: $\Gamma(x+1) = x\Gamma(x)$ ---
retrouvé sans une seule intégration par parties.

\textbf{21.} Avec $u = t^a$, $t = u^{1/a}$, $\dd t = \frac1a
u^{1/a - 1}\dd u$~:
\[
\int_0^\infty \eu^{-t^a}\dd t
= \frac1a\int_0^\infty u^{\frac1a - 1}\eu^{-u}\dd u
= \frac1a\,\Gamma\Bigl(\frac1a\Bigr)
= \Gamma\Bigl(1 + \frac1a\Bigr) .
\]
Quand $a \to +\infty$ (le long d'une suite quelconque)~: $\eu^{-t^a} \to 1$ pour
$0 < t < 1$, $\to \eu^{-1}$ en $t = 1$, $\to 0$ pour $t > 1$~;
pour $a \geq 2$ dominer par $\mathbf 1_{t \leq 1} +
\eu^{-t^2}\mathbf 1_{t > 1}$ ($t^a \geq t^2$ pour $t \geq 1$),
intégrable. Convergence dominée~: l'intégrale tend vers
$\int_0^1 1\,\dd t = 1$ --- comme il se doit, puisque $\Gamma(1 +
\frac1a) \to \Gamma(1) = 1$ par continuité.

\textbf{22.} Avec $u = t^n$, $\dd t = \frac1n u^{1/n - 1}\dd
u$~:
\[
\int_0^1 \frac{\dd t}{\sqrt{1-t^n}}
= \frac1n\int_0^1 u^{\frac1n-1}(1-u)^{-1/2}\dd u
= \frac1n\,B\Bigl(\frac1n, \frac12\Bigr) .
\]
$n = 1$~: $B\bigl(1,\frac12\bigr) = B\bigl(\frac12,1\bigr) = 2$,
correspondant à $\int_0^1\frac{\dd t}{\sqrt{1-t}} = 2$. $n = 2$~:
$\frac12 B\bigl(\frac12,\frac12\bigr) = \frac\pi2 = \arcsin 1$.
Pour $n = 4$ la valeur $\frac14 B\bigl(\frac14,\frac12\bigr)$ est
la constante de la lemniscate~: pas de forme close élémentaire.

\textbf{23.} En itérant l'équation fonctionnelle~:
\[
\frac{1}{\Gamma(x)}\int_0^\infty t^{x+k-1}\eu^{-t}\dd t
= \frac{\Gamma(x+k)}{\Gamma(x)}
= (x+k-1)(x+k-2)\cdots x ,
\]
la factorielle croissante à $k$ facteurs. En $x = 1$~:
$\Gamma(1+k)/\Gamma(1) = k!$, les moments de $\eu^{-t}$ de
l'\cref{exo:b2:integration:2}.

\textbf{24.} Substituons $t = \frac{1+s}2$ ($s \in
\intoo{-1}1$, $\dd t = \frac{\dd s}2$, $t(1-t) =
\frac{1-s^2}4$)~:
\[
B(x,x) = \int_{-1}^{1}\Bigl(\frac{1-s^2}{4}\Bigr)^{x-1}
\frac{\dd s}{2}
= 4^{1-x}\int_0^1 (1-s^2)^{x-1}\dd s
\]
(l'intégrande est pair). Puis $s = \sqrt v$ ($\dd s =
\frac{\dd v}{2\sqrt v}$)~:
\[
B(x,x) = \frac{4^{1-x}}{2}\int_0^1
v^{-1/2}(1-v)^{x-1}\dd v
= 2^{1-2x}\,B\Bigl(\frac12, x\Bigr) .
\]
Pour $2x \in \N^*$ tout argument en vue appartient à
$\frac12\N^*$, donc la formule d'Euler (question 14) s'applique aux deux
membres~:
\[
\frac{\Gamma(x)^2}{\Gamma(2x)}
= 2^{1-2x}\,
\frac{\Gamma\bigl(\frac12\bigr)\Gamma(x)}
{\Gamma\bigl(x+\frac12\bigr)}
\quad\Longleftrightarrow\quad
\Gamma(x)\,\Gamma\Bigl(x+\frac12\Bigr)
= 2^{1-2x}\sqrt\pi\;\Gamma(2x) .
\]
Vérification directe en $x = n$~: le membre de gauche est $(n-1)!\cdot
\frac{(2n)!\sqrt\pi}{4^n n!} = \frac{(2n)!\sqrt\pi}{4^n n}$,
le membre de droite $2\cdot4^{-n}\sqrt\pi\,(2n-1)! =
\frac{(2n)!\sqrt\pi}{4^n n}$~: égaux.

\textbf{25.} (i) L'intégration par parties a produit $B(x,y+1) =
\frac yx B(x+1,y)$, l'unique identité d'où découlent toute
relation de descente, les valeurs entières et demi-entières, et la
récurrence de Wallis. (ii) La convergence dominée a transformé les
intégrales élémentaires $\int_0^n(1-t/n)^n t^{x-1}$ en
$\Gamma(x)$ (formule de Gauss, question 17) et calculé la
limite $a \to \infty$ à la question 21. (iii) Du chapitre de comparaison
nous avons importé la constante d'Euler ($\ln n - H_n \to
-\gamma$, question 18) et l'équivalent du coefficient binomial central
(question 19) --- c.-à-d.\ la formule de Stirling déguisée. (iv)
La formule d'Euler $B(x,y) = \Gamma(x)\Gamma(y)/\Gamma(x+y)$ est
maintenant démontrée pour $y \in \N^*$ avec $x > 0$ arbitraire (question
10) et pour tous les couples demi-entiers (question 14)~; le cas
général $x, y > 0$ attend le calcul par intégrale double du
chapitre sur les intégrales multiples, qui factorise
$\Gamma(x)\Gamma(y)$ sur un quart de plan.
\end{solution}
